% !TEX root = userguide.tex

\def\headdate{4th July 2023}

\section{The convection-diffusion-reaction equation}

In this section some examples using the convection-diffusion-reaction
equation will be given. The examples will be split into a one-dimensional
description and a general one. The one-dimensional description is here
basically for educational reasons.

\subsection{Weak form of the steady convection-diffusion-reaction equation
 in 1D}
\label{sec:weakform1dconvdiff}

We consider the steady convection-diffusion-reaction equation on an interval
$\Omega=(a,b)$:
\begin{equation}
   u \deriv c x - \deriv{}x(\alpha\deriv{c}x) + \gamma c = f\text{ on }\Omega
\end{equation}
where $u$, $\alpha>0$ and $\gamma$ are constants and $f$ is a given function of
$x$.

The boundary conditions are
\begin{alignat*}{2}
   c&=c_a,\quad&&\text{ at }x=a\quad(\Gamma_D) \\
   -\alpha\deriv cx&=h_b \quad&&\text{ at }x=b\quad(\Gamma_N)
\end{alignat*}

The weak form is obtained by multiplying with a test function $v$ and
integrate:
\begin{equation}
   \big(v,u \deriv c x - \deriv{}x(\alpha\deriv{c}x) + \gamma c - f\big)\qquad
  \text{ for all }v
\label{eq:weak1}
\end{equation}
Partial integration of the diffusion term and inserting the boundary condition
at $x=b$ we get:
Find $c\in S$ such that
\begin{equation}
   (\deriv vx,\alpha\deriv cx)+(v,u\deriv cx)+(v,\gamma c)
        +v(b)h_b=(v,f)\qquad \text{for all }v\in V
\end{equation}
where $S$ and $V$ are appropriate spaces.

The weak form can be used to obtain an approximate solution using the
Galerkin \fem\ technique. For this we define approximating 
spaces $S_h$ and $V_h$, or in other words, we divide the region into elements:
\begin{equation}
\Omega=\cup_e\Omega_e, \quad\Omega_{e}\cap\Omega_{f}=\emptyset\quad
   \text{for }e\neq f
\end{equation}
and interpolate $c$ and $v$ by
polynomials on each element (approximation denoted by $c_h$ and $v_h$):
\begin{align*}
    c_h(x)&=\tsum_i c_i\phi_i(x)=\col\phi^T(x)\col c
\\
    v_h(x)&=\tsum_i v_i\phi_i(x)=\col\phi^T(x)\col v
\end{align*}
where $\col\phi(x)$ are global shape functions.  Substituting into the
weak form gives:
Find $c_h\in S_h$ such that
\begin{equation}
   (\deriv {v_h}x,\alpha\deriv {c_h}x)+(v_h,u\deriv {c_h}x)+
           (v_h,\gamma{c_h})+v_h(b)h_b
      = (v_h,f)\qquad \text{for all }v_h\in V_h
\end{equation}
or
\[
 \col v^T\mat K\col c=
        \col v^T\col f \quad\text{for all }\col v
\]
or
\[
 \mat K\col c= \col f
\]
where the `stiffness matrix' $\mat K$ and
the right-hand side $\col f$ are given by
\begin{alignat*}{2}
 \mat K&=(\deriv{\col\phi}{x},\alpha\deriv{\col\phi^T}{x})\qquad\text{  or}&
   \qquad K_{ij}&=(\deriv{\phi_i}{x}, \alpha\deriv{\phi_j}{x}) \\
      &\quad+(\col\phi,u\deriv{\col\phi^T}{x})  &
       &\quad+(\phi_i,u\deriv{\phi_j}{x})\\
      &\quad+(\col\phi,\gamma\col\phi^T)  &
       &\quad+(\phi_i,\gamma\phi_j)\\[2ex]
\col f&= (\col\phi,f)-h_b\col\phi(b) &
f_i&= (\phi_i,f)-h_b\phi_i(b)
\end{alignat*}
where the left column is in matrix notation and the right column in index
notation.

If convection dominates (P\'eclet number $\mathrm{Pe}=\dfrac{|u|L}{\alpha}\gg1$)
and the mesh is too coarse to resolve large gradients in the solution, the
Galerkin numerical solution can show `global wiggles'.
If mesh refinement is not an
option the so-called `upwind' schemes can help. The streamline upwind
Petrov-Galerkin (SUPG) \cite{Brooks82} is such a scheme and consists of
using a test function
$v+s\dfrac{\beta h}2 \deriv vx$ instead of $v$ in Eq.~\eqref{eq:weak1}:
\begin{equation}
   \big(v+s\dfrac{\beta h}2 \deriv vx,u \deriv c x - \deriv{}x(\alpha\deriv{c}x)
+ \gamma c - f\big)\qquad
  \text{ for all }v
\end{equation}
where $s$ is the sign\footnote{Strictly speaking, when derived from the 3D SUPG formulation, $s=u/|u|$, which in 1D  is the equal to sign of $u$.} of $u$, $h$ is the size of an element and the factor $\beta\geq0$ should be chosen
``optimally''. It should be noted that the additional terms must be interpreted
element wise, so that
\begin{itemize}
   \item $h$ and $\beta$ can vary from element to element,
   \item the multiplication of the test function with the second-order term can
be interpreted, i.e.
\[
   \big(v+s\dfrac{\beta h}2 \deriv vx, - \deriv{}x(\alpha\deriv{c}x)\big)=
   \big(v, - \deriv{}x(\alpha\deriv{c}x)\big) + \sum_e
   \big(s\dfrac{\beta h}2 \deriv vx, -\deriv{}x(\alpha\deriv{c}x)\big)
\]
The first term can be treated by partial integration like in the Galerkin
scheme giving the usual natural boundary conditions. Note, that the last term
is zero for linear elements and $\alpha$ is constant.
\end{itemize}
So the weak form for SUPG becomes:
Find $c\in S$ such that
\begin{multline}
   (\deriv vx,\alpha\deriv cx)-
 \sum_e \big(s\dfrac{\beta h}2 \deriv vx, \deriv{}x(\alpha\deriv{c}x)\big)
+(v+s\dfrac{\beta h}2 \deriv vx,u\deriv cx+\gamma c)
        +v(b)h_b\\
  =(v+s\dfrac{\beta h}2 \deriv vx,f)\qquad \text{for all }v\in V
\end{multline}
where $S$ and $V$ are appropriate spaces.

The parameter $\beta$ depends on the so-called mesh P\'eclet number 
$\mathrm{Pe}_h=\dfrac{|u|h}{2\alpha}$.
For accuracy, $\beta$ should become small 
and tend to zero for $\mathrm{Pe}_h<1$ (tend to the Galerkin method). For
example 
\[
    \beta=\coth(\mathrm{Pe}_h)-\frac1{\mathrm{Pe}_h}
\]
which gives exact nodal values if $f$ is a constant, or
\[
   \beta=
\begin{cases}
    0  & \text{ for }\mathrm{Pe}_h\leq1 \\
    1-\dfrac1{\mathrm{Pe}_h} & \text{ for }\mathrm{Pe}_h>1
\end{cases}
\]
which keeps the `artificial' mesh P\'eclet number equal to 1 for
$\mathrm{Pe}_h>1$ in case of the artificial diffusion method (SU, see below).

The streamline upwind (SU) method is an inconsistent method and is only
mentioned here for comparison. It is constructed by applying
the upwind biased test function to
the convection term only:
\begin{equation}
   (v+s\dfrac{\beta h}2 \deriv vx,u \deriv c x) +\big(v, -
\deriv{}x(\alpha\deriv{c}x) + \gamma c - f\big)\qquad
  \text{ for all }v
\end{equation}
It should not be used, because the improper weighting leads to large errors.
Note, however that for $f=0$, $\gamma=0$ and and in case of
linear elements the SUPG and SU
method are identical.

\subsection{Examples of the steady 1D convection-diffusion-reaction}
\label{sec:examples1dconvdiff}

These examples solves the 1D convection-diffusion-reaction equation on $(0,L)$
with constant coefficients for various source terms. 
In the first example (\texttt{convdiff1.f90}), the three methods Galerkin,
SUPG and SU, using standard \fem, can be used for comparison and generating textbook examples.
In the second example (\texttt{convdiff2.f90}), the Galerkin method is used together
with spectral elements.

The example files can be obtained by typing at the
command line:
\medskip
\begin{quote}
   \texttt{getexample convection\_diffusion\_1d}
\end{quote}
The examples are self-contained and do not use add-on element routines.
Therefore it can be used to study a simple complete example without bothering
about complicated add-ons.

\subsection{Weak form of the convection-diffusion-reaction
equation. Galerkin \fem}
\label{sec:diffusionweak}

We consider the convection-diffusion-reaction equation in a region $\Omega$:
\begin{equation}
    \gamma(\pderiv{c}t + \vek u\cdot\nabla c) -\nabla \cdot( \alpha \nabla c)
+ \beta c = f\text{ on }\Omega
 \label{eq:convectiondiffusion}
\end{equation}
where the coefficient $\alpha$ is either a scalar or a tensor. All the
coefficients ($\alpha$, $\beta$, $\gamma$, $\vek u$ and $f$) may vary in
space and time and may also depend on other quantities.
For $\gamma=\beta=0$, this equation reduces to
the (steady) diffusion equation.
If in addition $\alpha=1$ the equation reduces to
the Poisson equation (see Sec.~\ref{sec:poissonweak}).
The boundary conditions are
\begin{gather}
    c = c_D\text{ on }\Gamma_D\\
    -\vec n\cdot(\alpha\nabla c)=h_{\text{N}}\text{ on }\Gamma_{\text{N}}
\end{gather}
where the boundary $\Gamma=\Gamma_D\cup\Gamma_{\text{N}}$ has been split into a
Dirichlet and a Neumann part.

The weak form can be obtained by multiplying Eq.~\eqref{eq:convectiondiffusion}
with a test function $v$:
\[
   \big(v,\gamma(\pderiv{c}t + \vek u\cdot\nabla c)
     -\nabla\cdot( \alpha \nabla c)+ \beta c-f\big)=0\qquad\text{for all }v
\]
where the standard inner product on $L^2(\Omega)$ is
\[
    (a,b)=\int_\Omega ab \D x
\]
With
\[
   v\nabla\cdot\big(\alpha\nabla c\big)=\nabla\cdot\big(v\alpha\nabla c)\big)
-\nabla v\cdot\alpha\nabla c
\]
and the divergence theorem (Gauss):
\[
    \int_\Omega\nabla\cdot\vek a\D x=\int_\Gamma\vek n\cdot\vek a\D s
\]
we get the weak form: Find $c\in S$ such that
\begin{equation}
   (v,\gamma\pderiv{c}t)+(v,\gamma \vek u\cdot\nabla c)
   +(\nabla v,\alpha\nabla c) + (v,\beta c) +
             (v,h_{\text{N}})_{\Gamma_{\text{N}}}=(v,f)
   \label{eq:weakformdiffusion}
\end{equation}
for all $v\in V$, where $S$ and $V$ are suitable functional spaces for $c$ and
$v$, respectively. Notes:
\begin{itemize}
  \itemsep -1ex
   \item $v=0$ on $\Gamma_D$, $c=c_D$ on $\Gamma_D$
   \item we have silently introduced:
\begin{align*}
   (\vek a, \vek b)&=\int_\Omega\vek a\cdot\vek b\D x\\
   (a,b)_{\Gamma}&=\int_\Gamma a b\D s
\end{align*}
\end{itemize}

The weak form can be used to obtain an approximate solution using the Galerkin \fem\
technique. For this we define approximating 
spaces $S_h$ and $V_h$, or in other words, we divide the region into elements:
\[
\Omega=\cup_e\Omega_e, \quad\Omega_{e}\cap\Omega_{f}=\emptyset\quad
   \text{for }e\neq f
\]
and interpolate $c$ and $v$ by
polynomials on each element (approximation denoted by $c_h$ and $v_h$):
\begin{align*}
    c_h(\vek x,t)&=\tsum_i c_i(t)\phi_i(\vek x)=\col\phi^T(\vek x)\col c(t)
\\
    v_h(\vek x,t)&=\tsum_i v_i(t)\phi_i(\vek x)=\col\phi^T(\vek x)\col v(t)
\end{align*}
where $\col\phi(\vek x)$ are global shape functions.  Substituting into the
weak form leads to
\[
 \col v^T(\mat M\dot{\col c}+\mat K\col c)=
        \col v^T\col f \quad\text{for all }\col v
\]
or
\begin{equation}
 \mat M\dot{\col c}+\mat K\col c= \col f
 \label{eq:convdiffdiscr}
\end{equation}
where the `mass matrix' $\mat M$,
the `stiffness matrix' $\mat K$ and
the right-hand side $\col f$ are given by
\begin{align*}
 \mat M&=(\col\phi,\gamma\col\phi^T)\\
\mat K&=(\col\phi,\gamma\vek u\cdot\nabla\col\phi^T)+
(\nabla{\col\phi},\alpha\nabla{\col\phi^T})+(\col\phi,\beta\col\phi^T)\\
\col f&= (\col\phi,f)-(\col\phi,h_{\text{N}})_{\Gamma_{\text{N}}}
\end{align*}
or in index notation
\begin{align*}
M_{ij}&=(\phi_i, \gamma\phi_j) \\
K_{ij}&=(\phi_i,\gamma\vek u\cdot\nabla\phi_j)
+(\nabla{\phi_i}, \alpha\nabla{\phi_j})+(\phi_i, \beta\phi_j) \\
f_i&= (\phi_i,f)-(\phi_i,h_{\text{N}})_{\Gamma_{\text{N}}}
\end{align*}

The element routines have been designed such that many standard time
discretization schemes for \eqref{eq:convdiffdiscr}
can be programmed in an easy way.
A first class of time discretization methods is based on backwards differencing.
The most simple method is the implicit Euler (IE) or backward Euler 
scheme:
\[
 \mat M\frac{\col c^{n+1}-\col c^n}{\Delta t}+\mat K\col c^{n+1}= \col f^{n+1}
\]
where the coefficient matrices $\mat M$ and $\mat K$ can be evaluated at either
the new $t^{n+1}$ or the old $t^n$ time.
It is first-order accurate ($\mathcal{O}(\Delta t)$), but very stable. A
natural extension is using second-order backwards
differencing (BDF2) (the second-order Gear scheme, see \cite{Hulsen96} for
further references):
\[
 \mat M^{n+1}\frac{\frac32\col c^{n+1}-2\col c^n+\frac12\col c^{n-1}}{\Delta t}+
\mat K^{n+1}\col c^{n+1}= \col f^{n+1}
\]
This is a multi-step scheme. It needs to be started, for example by the implicit Euler method.

A simple one-step second-order method is based on the trapezoidal rule. 
The fully implicit trapezoidal rule is also called Crank-Nicolson (CN) or second-order Adams-Moulton (AM2):
\[
 \mat M\frac{\col c^{n+1}-\col{c}^n}{\Delta t}+
\frac12(\mat K^{n+1}\col c^{n+1}+
\mat K^{n}\col c^{n})=\frac12( \col f^{n+1}+ \col f^{n})
\]
where the method has been defined using a constant mass matrix $\mat M$.
Although the CN-method has a small error coefficient (smaller than BDF2),
in practice large
oscillating errors can be present due to the bad damping properties
(amplification factor goes to -1 for large steps). A possible way of avoiding
these oscillations is starting with either an implicit Euler scheme
or a ``theta-scheme''
\[
 \mat M\frac{\col c^{n+1}-\col{c}^n}{\Delta t}+
\theta\mat K^{n+1}\col c^{n+1}+
(1-\theta)\mat K^{n}\col c^{n}=\theta\col f^{n+1}+ (1-\theta) \col f^{n}
\]
with $\theta\in[1,\frac12)$,
for several time steps to damp the oscillations.

\subsection{A steady diffusion problem on a unit square with Dirichlet
            and Neumann boundary conditions. Varying diffusion coefficient and
            (optionally) a varying Gauss rule}

\label{sec:diffusion1}

We consider the steady
diffusion problem\footnote{The example, and other examples derived from it,
is rather artificial. Its purpose is to show the use of the 
\texttt{convection\_diffusion\_m} module.}.
 on a unit square, similar to the Poisson
problem in Sec.~\ref{sec:poisson2}. We use the same function for $f$ on the
square domain and $\plderiv{c}{n}$ on the
curve \texttt{C2} as in the Poisson problem. However, now the diffusion
coefficient is not constant but given by the function
\begin{equation}
     \alpha=\tfrac32+\cos(\pi x)\cos(\pi y)
\end{equation}
We don't know the exact solution for $c$, but that is not important here.

The sources are
in the files \texttt{diffusion1.f90}, \texttt{diffusion2.f90} and
\texttt{diffusion3.f90}.
 and can be obtained by typing at the
command line:
\begin{quote}
   \texttt{getexample diffusion}
\end{quote}
In \texttt{diffusion1.f90} the problem can be solved by
either directly calling a function routine for $\alpha$ in the Gauss points
 on element level or by filling a vector (defined in the nodes) containing
the value of $\alpha$ using the same function. In the latter case the Gauss
point values of $\alpha$ are determined by using the same shape function as
used for $u$.
In \texttt{diffusion2.f90} the problem is solved by filling a vector defined
per element of type \texttt{elvector\_t} with the Gauss point values of
$\alpha$.
In \texttt{diffusion3.f90} the problem is also
solved by filling a vector defined
per element of type \texttt{elvector\_t} with the Gauss point values of
$\alpha$. However, in addition, here it is shown how to use a varying
Gauss integration rule, i.e.\ the rule can be different for each element.

\subsection{A time-dependent diffusion problem on a unit square with Dirichlet
            and Neumann boundary conditions. Varying diffusion coefficient.
            Second-order Gear time-discretization}

\label{sec:diffusion2}

We consider the same 
diffusion problem on a unit square of Sec.~\ref{sec:diffusion1} as given by
the source code \texttt{diffusion1.f90}, however
now a time-dependent problem is solved with $\gamma=1$.
The source is \texttt{diffusion4.f90}.
The time-discretization is second-order Gear.

\subsection{A steady convection-diffusion problem on a unit square with
            Dirichlet and Neumann boundary conditions.
            Varying diffusion and velocity coefficient.
            Scalar to tensor diffusion coefficient.}

\label{sec:convectiondiffusion1}

We consider the same 
diffusion problem on a unit square of Sec.~\ref{sec:diffusion1} as given by
the source code \texttt{diffusion1.f90}, however
now a problem is extended with a convection term with varying velocity
coefficient.

The sources are
in the files \texttt{convection\_diffusion1.f90}
and \texttt{convection\_diffusion2.f90} 
 and can be obtained by typing at the
command line:
\begin{quote}
   \texttt{getexample convection\_diffusion}
\end{quote}
In \texttt{convection\_diffusion1.f90} the diffusion coefficient is a scalar
and in \texttt{convection\_diffusion2.f90} it is a tensor.

\subsection{A time-dependent convection-diffusion problem on a unit square with
            Dirichlet and Neumann boundary conditions.
            Varying diffusion and velocity coefficient.
            Varying $\beta$ coefficient.
            Second-order Gear (BDF2) or Crank-Nicolson (CN)
            time-discretization}

\label{sec:convectiondiffusion2}

We consider the same 
convection-diffusion problem on a unit square
of Sec.~\ref{sec:convectiondiffusion1} as given by
the source code \texttt{convection\_diffusion1.f90}, however
now a time-dependent problem is solved with $\gamma=1$.

The sources are
in the files \texttt{convection\_diffusion3.f90} (BDF2),
\texttt{convection\_diffusion4.f90} (BDF2, $\beta$-coefficient),
\texttt{convection\_diffusion5.f90} (CN, $\beta$-coefficient) and
\texttt{convection\_diffusion5.f90} (CN, $\beta$-coefficient, tensor $\alpha$) 
and can be obtained by typing at the
command line:
\begin{quote}
   \texttt{getexample convection\_diffusion}
\end{quote}
The examples using CN are far from optimal and show oscillations in time.

\subsection{Steady 1D convection-diffusion equation using spectral elements}
\label{sec:convdiff1D}
This example is similar to the example \texttt{convdiff2.f90}
in Sec.~\ref{sec:examples1dconvdiff}, but now using the standard elements in the add-on
\texttt{convection\_diffusion}.
The source is in the file \texttt{convection\_diffusion7.f90}
and can be obtained by typing at the
command line:
\begin{quote}
   \texttt{getexample convection\_diffusion}
\end{quote}

\subsection{A time-dependent diffusion problem with \texttt{vtk}-output}

This example solves a time-dependent temperature problem in a cylindrical rod.
This problem solves the example in Sec.~3.2.2 of Dantzig \& Tucker \cite{Dantzig2001}:
A solid cylindrical rod of radius $R$ and length $L$, at an initial
temperature $T_1$, is quenched in a fluid at a lower temperature $T_0$.
Second-order implicit Gear discretization. First step: Euler implicit.
Half the domain is solved, due to symmetry.
Output of \texttt{vtk}-files as a function of time.

The source is in the file \texttt{diffusion9.f90}
and can be obtained by typing at the
command line:
\begin{quote}
   \texttt{getexample diffusion}
\end{quote}

\subsection{Stabilization of the convection-diffusion-reaction equation using SUPG}
\label{sec:convdiffsupg}

We consider the following generic convection-diffusion-reaction equation for 
$c(\vek x, t)$
\begin{equation}
  \gamma \left( \pderiv{c}{t} + \vek u \cdot \nabla c \right) - \nabla \cdot 
  (\alpha \nabla c) + \beta c = f, \label{eq:conv_diff}
\end{equation}
where $\gamma$, $\alpha$ and $\beta$ are constant coefficients and $f$ is a 
source 
term, which can be position and time dependent (i.e.~$f=f(\vek x,t)$). The 
weak form of Eq.~\eqref{eq:conv_diff} is given by: find $c$ such that
\begin{equation}
  \gamma (v, \pderiv{c}{t}+\vek u \cdot \nabla c) + \alpha (\nabla v, \nabla 
  c)+\beta(v,c)-\alpha (v, \vek n \cdot \nabla c)_\Gamma = (v, f ) 
  \label{eq:conv_diff_weak}
\end{equation}
for all admissible test functions $v$. A flux boundary condition on 
$\Gamma$ (with unit normal $\vek n$) can be substituted directly into equation 
\eqref{eq:conv_diff_weak}, but for now we will assume this term to be zero.
To integrate \eqref{eq:conv_diff_weak} in time, we assume that the velocity 
$\vek u$ and 
source term $f$ at some time $t_{n+1}$ are known. Using a 
first-order backward
Euler scheme for Eq.~\eqref{eq:conv_diff_weak} yields:
\begin{equation}
  \gamma (v, \frac{c_{n+1}-c_n}{\Delta t}+\vek u \cdot \nabla c_{n+1}) + \alpha 
  (\nabla v, \nabla c_{n+1}) +\beta(v,c_{n+1})= (v, f ) 
  \label{eq:conv_diff_euler}
\end{equation}
whereas a second-order Gear scheme (BDF2) yields
\begin{equation}
  \gamma (v, \frac{\frac{3}{2}c_{n+1}-2c_n+\frac{1}{2}c_{n-1}}{\Delta t}+\vek u 
  \cdot \nabla c_{n+1}) + \alpha (\nabla v, \nabla c_{n+1}) +\beta(v,c_{n+1})= (v, f ). 
  \label{eq:conv_diff_gear}
\end{equation}

The ratio between the convection and diffusion terms is expressed in the 
P\'{e}clet number $\mathrm{Pe}$. For high $\mathrm{Pe}$ (transport is convection-dominated), 
which 
might lead to ``wiggles'' if the gradients in the solution are not properly 
resolved. Streamline upwind Petrov-Galerkin (SUPG) can help in reducing these wiggles, 
as is explained in Section \ref{sec:weakform1dconvdiff}. Using SUPG, 
Eq.~\eqref{eq:conv_diff_weak} is written as: find $c$ such that
\begin{multline}
  \gamma (v+\tau \vek u \cdot \nabla v, \pderiv{c}{t} + \vek u \cdot 
  \nabla c + \beta c)  + \alpha 
  (\nabla v, \nabla c) - \sum_\text{e} (\tau \vek u \cdot \nabla v,\nabla 
  \cdot (\alpha \nabla c)) \\
  = (v+\tau \vek u \cdot \nabla v, f ) 
  \label{eq:conv_diff_supg}
\end{multline}
for all admissible test functions $s$. Here $\tau = \beta_\textsc{supg} h / (2U)$, where $h$ 
is the element size in the direction 
of the velocity and $U$ is the (local) characteristic velocity. Details on how to 
compute $h$ and $U$ can be found in Appendix \ref{chap:tausupg}. Note, that the 
last term on the left-hand side still contains second-order derivatives of $c$, 
but these will be zero if linear base functions (on a line, triangle or tetrahedron) are used. 
For other base functions, this term will simply be ignored. The value of $\beta_\textsc{supg}$ 
varies between zero and one, and can be made dependent on the local (mesh) 
P\'{e}clet number, which is defined as:
\begin{equation}
  \mathrm{Pe}_h = \frac{\gamma U h}{2\alpha}.\label{eq:Peh}
\end{equation}
The value of $\beta_\textsc{supg}$ should approach zero for small 
$\mathrm{Pe}_h$. For linear base functions $\beta_\textsc{supg}$ should approach one for large values of 
$\mathrm{Pe}_h$. 

\subsection{A 2D unsteady convection problem with SUPG stabilization}
An example
of SUPG in a 2D unsteady convection problem is available in the example
\texttt{convection\_diffusion\_supg1.f90} and can be obtained  
by typing at the command line:
\begin{quote}
   \texttt{getexample convection\_diffusion\_supg}
\end{quote}

\subsection{Multiple coupled convection equations with SUPG stabilization}

Problems in crystallization can be described by multiple (coupled) equations of the form of 
Eq.~\eqref{eq:conv_diff} but without the diffusion term, e.g.~:
\begin{equation}
  \pderiv{c_1}{t} + \vek u \cdot \nabla c_1  =  0, \qquad
  \pderiv{c_2}{t} + \vek u \cdot \nabla c_2 =  c_1.
  \label{eq:conv_coupled}
\end{equation}  
The approach taken here is to solve these equations decoupled, i.e.~first we 
solve for $c_1$ and then we solve for $c_2$. Dependencies can be taken into 
account by adding the updated solution to 
the right-hand side of the consequent equations (e.g.~$(c_1)_{n+1}$ is solved 
first and is then added to the right hand side to find $(c_2)_{n+1}$). This 
approach has 
the advantage that updated solutions are used in the solving procedure, at the 
expensive of additional assembly steps of the right-hand side. A second approach
is to assemble the right-hand side using a prediction for $c_1$, 
i.e.~$(\hat{c}_1)_{n+1}= (\hat{c}_1)_{n}$ for first-order, or 
$(\hat{c}_1)_{n+1}= 2(\hat{c}_1)_{n}-(\hat{c}_1)_{n-1}$ for second-order time 
intergration). With this approach, the system can be assembled in one go, saving
time in the assembly.  Note that if $\vek u$ is constant, the system matrix for 
these equations does not change, thus the LU-decomposition only has be 
computed once and can be re-used. If $\vek u$ varies in time, the 
LU-decomposition can still be re-used within a time step. Note, that the 
Eq.~\eqref{eq:conv_coupled} is hyperbolic, and SUPG can be used for 
stabilization. In the examples \texttt{convection\_diffusion\_supg2.f90} and  
\texttt{convection\_diffusion\_supg3.f90}, the two equations as shown in 
Eq.~\eqref{eq:conv_coupled} are solved using a manufactured solution. In the
former example, the updated solution is used to fill the right-hand side, 
whereas in the latter example, the right hand side is filled using a prediction
for $c_1$. The examples can be obtained by typing at the command line:
\begin{quote}
   \texttt{getexample convection\_diffusion\_supg}
\end{quote}

